\section{Related Work}

\paragraph{Low-Precision training.}
Low-precision training has a long history, including early work on reduced-precision and integer training for smaller neural networks \citep{courbariaux2015binaryconnect,banner2018scalable,yang2019training}. For Transformer models, 8-bit training has become increasingly practical through INT8 methods such as SwitchBack and JetFire, FP8 training systems, and 8-bit LLM fine-tuning methods such as HALO \citep{wortsman2023stablelowprecisiontraininglargescale, xi2024jetfireefficientaccuratetransformer, micikevicius2022fp8, deepseekai2025deepseekv3technicalreport, mishra2025recipespretrainingllmsmxfp8, ashkboos2025halohadamardassistedlowerprecisionoptimization}. These methods highlighted several recurring themes: the need for fine-grained scaling, special treatment of outlier values, and careful handling of the backward pass.

\paragraph{Training in FP4 formats.}
The introduction of hardware-supported FP4 microscaling formats has renewed interest in end-to-end 4-bit training. Early FP4 and INT4 work demonstrated that 4-bit arithmetic could be incorporated into training \citep{xi2023trainingtransformers4bitintegers, chmiel2024accurateneuraltraining4bit, wang2025optimizinglargelanguagemodel, zhou2025efficientpretrainingexploringfp4}. More recent work targets hardware-supported MXFP4 and NVFP4 training. \citet{tseng2025mxfp4} studied backward-pass quantization to FP4 using Hadamard rotations and stochastic rounding. The TetraJet series combined one-dimensional weight quantization with re-quantization in the backward pass and stochastic rounding, and developed methods for oscillation suppression; TetraJet-v2 extended the approach to NVFP4 LLM training, adding double-block quantization, rotations to both backward GEMMs, and an outlier control mechanism \citep{chen2025oscillationreducedmxfp4trainingvision, chen2025tetrajet}. \citet{chmiel2025fp4all} studied fully quantized FP4 training with NVFP4-style scaling and identified stochastic rounding as important for the backward and update passes. \citet{nvidia2025pretraining} introduced a large-scale NVFP4 recipe based on Random Hadamard transforms, two-dimensional weight quantization, stochastic rounding, and selective high-precision layers. Four Over Six \citep{cook2025fouroversix} improved NVFP4 quantization accuracy by adaptively choosing between two scale factors.


\begin{figure*}[h!]
\centering
  \includegraphics[width=\textwidth]{figures/icml_workshop_fig1.png}
  \caption{Left: Relative train loss difference between each NVFP4 recipe and BF16 (lower is better), trained with \nvrecipe and \newrecipe recipes to 1T tokens on Nemotron 3 Nano. Right: Cumulative recipe updates, each adding one \newrecipe component, trained to 400B tokens on Nemotron 3 Nano. All recipes keep the last 15\% of layers in BF16. Recipe updates are described in Table~\ref{recipe_table}.}
  \label{fig:recipe_stacking}
\end{figure*}

\paragraph{Quartet and unbiased gradient estimation.}
Quartet \citep{castro2025quartet} studied MXFP4 training through the lens of ``low-precision scaling laws,'' linking forward compression error to parameter efficiency and backward-pass bias to data efficiency. Historically, methods addressing bias in quantized gradient estimates have generally relied on elementwise stochastic rounding \citep{chmiel2024accurateneuraltraining4bit, tseng2025mxfp4, chen2025oscillationreducedmxfp4trainingvision, chmiel2025fp4all, nvidia2025pretraining, chen2025tetrajet}. Quartet II \citep{panferov2026quartetii} extended the Quartet line of work to NVFP4 and introduced MS-EDEN, which instead combines random rotations with a per-block corrective rescaling to produce lower-error unbiased estimates than standard FP4 stochastic rounding. Our work is closest in spirit to these studies, but asks which ingredients of existing FP4 training recipes matter most, and how they should be combined, when evaluated under a common implementation and training budget.

